\documentclass{article}
\usepackage[UTF8]{ctex}
\usepackage{amsmath,amssymb,amsfonts}
\usepackage{enumitem}
\usepackage{geometry}
\geometry{a4paper, margin=1in}
% === 错题本排版优化（不参与题目标识） ===
\setlength{\parindent}{0pt}
\setlength{\parskip}{0.35em plus 0.10em minus 0.08em}
\setlength{\emergencystretch}{3em}
\setlist{itemsep=0.25em, topsep=0.40em, parsep=0pt}
\allowdisplaybreaks[2]
% === 错题本排版优化结束 ===
\begin{document}
\section*{东北师范大学2026年数学分析真题}
\section*{80．东北师范大学2026年数学分析试卷（回忆版）}

\subsection{题目1}
1．求极限 \(\lim _{n \rightarrow \infty}\left(2026 \cos ^{2} n+\sin ^{2} n\right)^{\frac{1}{n}}\) ．

\subsection{题目2}
2 ．求极限 \(\lim _{x \rightarrow 0}\left(\frac{1^{x}+2^{x}+\cdots+2026^{x}}{2026}\right)^{\frac{1}{\sin x}}\) ．

\subsection{题目3}
3．计算重积分 \(\iint_{D}\left(x^{2}+y^{2}\right) \mathrm{d} x \mathrm{~d} y\) ，其中 \(D\) 由椭圆 \(9 x^{2}+y^{2}=1\) 所围成。

\subsection{题目4}
4．想不起来了。

\subsection{题目5}
5．计算定积分 \(\int_{-1}^{1} \frac{t^{2027}-t^{2025}+t+1}{\cos ^{2} t} \mathrm{~d} t\) ．

二．（15 分）计算曲线积分 \(\oint_{L} \frac{x \mathrm{~d} y-y \mathrm{~d} x}{16 x^{2}+y^{2}}\) ，其中 \(\Gamma\) 为圆周 \((x-2)^{2}+y^{2}=9\) ，取逆时针方向．

三．（15 分）设函数 \(u(x, y)\) 由方程组

\[
\left\{\begin{array}{l}
y z+2 x z+x y=u \\
x^{2}+y^{2}+z^{2}=1
\end{array}\right.
\]

所确定，求 \(\frac{\partial u}{\partial x}\) 及 \(\frac{\partial^{2} u}{\partial y \partial x}\) ．

四．（15 分）设 \(k(x, y) \in C([0,1] \times[0,1]), f(x)\) 在 \([0,1]\) 上可积，且

\[
h(x)=\int_{0}^{x} f(y) k(x, y) \mathrm{d} y
\]

证明：\(h(x)\) 在 \([0,1]\) 上连续．

五．（15 分）求由曲线 \(y=2 x^{2}\) 与 \(y=x\) 所围成的平面图形的面积，并求

\begin{itemize}
  \item[（1）] 该图形绕 \(x\) 轴旋转所得的旋转体体积。
  \item[（2）] 该图形绕 \(y\) 轴旋转所得的旋转体体积。
\end{itemize}

六．（15 分）求函数 \(f(x, y)=x^{2}+y^{2}-6 x+8 y\) 在区域 \(D: x^{2}+y^{2} \leq 25\) 上的最大值与最小值．

七．（20 分）设函数列 \(f_{n}(x)=x^{2025 n}-x^{2026 n}\) ．

\begin{itemize}
  \item[（1）] 证明函数列 \(\left\{f_{n}(x)\right\}\) 在 \(\left[0, \frac{2025}{2026}\right]\) 上一致收敛．
  \item[（2）] 说明函数列 \(\left\{f_{n}(x)\right\}\) 在 \([0,1]\) 上不一致收敛．
\end{itemize}

八．（15 分）设 \(f(x)\) 在 \([0,1]\) 上连续，在 \((0,1)\) 上可导，且

\[
f^{2}(1)-f^{2}(0)=2026
\]

证明：存在 \(\xi \in(0,1)\) ，使得 \(f^{\prime}(\xi) f(\xi)=2026 \xi\) ．

九．（10 分）设正数列 \(\left\{a_{n}\right\}\) 满足 \(a_{n+1} \leq a_{n}+\frac{1}{2026^{n}}\) ．

\begin{itemize}
  \item[（1）] 证明数列 \(\left\{a_{n}\right\}\) 有界．
  \item[（2）] 证明数列 \(\left\{a_{n}\right\}\) 收敛．
\end{itemize}
\end{document}
